\chapter{PENDAHULUAN}

\section{Latar Belakang}
Kehidupan akademik mahasiswa menuntut pengelolaan banyak hal sekaligus:
tugas kuliah dengan tenggat yang berbeda-beda, jadwal perkuliahan, email
dari dosen dan kampus, serta koordinasi kerja kelompok. Informasi tersebut
tersebar di berbagai kanal --- email kampus, grup pesan instan, sistem
pembelajaran daring, dan pengumuman lisan di kelas --- sehingga mahasiswa
sering melewatkan tenggat atau kesulitan menentukan prioritas pengerjaan.

Kemampuan manajemen waktu terbukti berpengaruh positif terhadap prestasi
akademik mahasiswa \parencite{kamilatunnisa2025}. Sebaliknya, kegagalan
regulasi diri seperti prokrastinasi akademik berhubungan erat dengan
menurunnya kinerja dan kesejahteraan mahasiswa \parencite{fadila2024}.
Sayangnya, aplikasi manajemen
tugas yang umum digunakan (daftar tugas dan kalender generik) masih
menuntut pengguna memasukkan seluruh data secara manual, padahal sumber
informasi tugas terbesar mahasiswa justru berada di email kampus dan
dokumen perkuliahan yang tidak terstruktur.

Perkembangan kecerdasan buatan, khususnya model bahasa besar
(\emph{large language model}/LLM), membuka peluang untuk mengotomatiskan
pengelolaan informasi akademik dan terbukti bermanfaat sebagai asisten
pembelajaran digital bagi mahasiswa
\parencite{johan2026, hidayah2025}. Email kampus dapat dianalisis untuk mendeteksi tugas
secara otomatis, tugas dapat dipecah menjadi subtask dengan estimasi waktu,
dan seluruh tenggat dapat disajikan dalam satu dasbor terpadu.

Berdasarkan permasalahan tersebut, tugas akhir ini mengusulkan pengembangan
\textbf{Welya}, asisten akademik berbasis AI dalam bentuk produk digital
(aplikasi web dan desktop) yang membantu mahasiswa mengelola tugas, jadwal, email
kampus, dan kerja kelompok dalam satu tempat. Welya terintegrasi dengan
Gmail, Google Tasks, dan Google Calendar sehingga tugas terdeteksi
otomatis, tersusun rapi, dan mudah dipantau melalui satu dasbor.

\section{Rumusan Masalah}
\begin{enumerate}
  \item Bagaimana merancang dan membangun asisten akademik berbasis AI
    yang mampu mendeteksi tugas secara otomatis dari email kampus?
  \item Bagaimana mengintegrasikan aplikasi dengan layanan Gmail, Google
    Tasks, dan Google Calendar agar tugas dan jadwal mahasiswa
    terkelola dalam satu dasbor?
  \item Bagaimana tingkat kebermanfaatan dan kemudahan penggunaan
    aplikasi bagi mahasiswa?
\end{enumerate}

\section{Tujuan}
\begin{enumerate}
  \item Menghasilkan produk digital Welya berupa aplikasi web dan desktop asisten
    akademik berbasis AI yang siap diuji coba.
  \item Mengimplementasikan integrasi Gmail, Google Tasks, dan Google
    Calendar untuk deteksi tugas dan sinkronisasi jadwal otomatis.
  \item Mengukur kebermanfaatan dan kemudahan penggunaan aplikasi melalui
    uji coba kepada mahasiswa.
\end{enumerate}

\section{Manfaat}
\begin{enumerate}
  \item Bagi mahasiswa: membantu mengatur tugas, jadwal, dan email kampus
    sehingga kuliah lebih teratur dan tenggat tidak terlewat.
  \item Bagi institusi: menjadi contoh penerapan teknologi AI dalam
    mendukung produktivitas akademik mahasiswa.
  \item Bagi penulis: menerapkan ilmu rekayasa perangkat lunak dan
    kecerdasan buatan dalam produk nyata yang berpotensi memperoleh
    HKI.
\end{enumerate}

\section{Batasan Masalah}
\begin{enumerate}
  \item Aplikasi tersedia sebagai aplikasi web dan aplikasi desktop
    serta menargetkan pengguna mahasiswa.
  \item Integrasi layanan pihak ketiga terbatas pada Gmail, Google Tasks,
    dan Google Calendar.
  \item Deteksi tugas otomatis berfokus pada email kampus berbahasa
    Indonesia.
  \item Uji coba dilakukan pada lingkup terbatas mahasiswa sebagai
    pengguna awal.
\end{enumerate}
